\chapter{Limitationen und Ausblick}
\label{sec:limitationen}

\section{Limitationen}

Die Arbeit weist mehrere Einschränkungen auf. Die Anzahl der befragten Expertinnen und Experten ist begrenzt. Dadurch sind die Ergebnisse nicht uneingeschränkt verallgemeinerbar. Zudem basieren automatisierte Evaluationsmetriken wie RAGAS auf LLM-basierten Bewertungsmodellen, die modellabhängig sind und keine objektive Ground Truth darstellen. Die Ergebnisse spiegeln daher die Eigenschaften der eingesetzten Modelle wider und können in ihrer Reproduzierbarkeit variieren.

Aufgrund technischer Schwierigkeiten mit der virtuellen Maschinenumgebung konnten einzelne quantitative Messungen nicht vollständig erfasst und ausgewertet werden. Die Datenerhebung wird weiterhin fortgeführt; eine umfassende Auswertung ist im Rahmen einer geplanten Veröffentlichung im März vorgesehen.

Die Evaluation erfolgte im spezifischen Anwendungskontext des KI-Campus innerhalb eines Learning-Management-Systems. Die Ergebnisse sind daher primär für diese Systemarchitektur, Datenbasis und Zielgruppe gültig. Auch die beobachtete Unsicherheit bei mehrsprachigen, insbesondere englischen Anfragen, wurde nicht systematisch isoliert untersucht.

Die tatsächliche Lernwirkung wurde nicht experimentell gemessen, sondern auf Basis qualitativer Einschätzungen bewertet. Ebenso ist die Kostenanalyse modell- und anbieterspezifisch, sodass Veränderungen in Preisstrukturen oder Modellversionen zu abweichenden Bewertungen führen können.

\section{Ausblick}

Zukünftige Forschung sollte die tatsächliche Lernwirkung experimentell untersuchen, beispielsweise durch kontrollierte oder longitudinale Studiendesigns in realen Kursumgebungen des KI-Campus. Ebenso erscheint eine systematische Analyse unterschiedlicher Routing-Strategien unter variierenden Nutzungsszenarien sinnvoll, um kontextspezifische Optimierungsstrategien zu identifizieren.

Darüber hinaus sollte der Einsatz spezialisierter Reranker-Modelle hinsichtlich ihres Kosten-Nutzen-Verhältnisses empirisch geprüft werden. Auch eine gezielte architektonische Verbesserung der Mehrsprachigkeit stellt ein relevantes Entwicklungsfeld dar. Schließlich bleibt die Weiterentwicklung transparenter Evaluationsmethoden sowie die Berücksichtigung ethischer und regulatorischer Anforderungen ein zentrales Forschungsfeld.

Insgesamt liefert die Arbeit eine empirisch fundierte Grundlage für die Weiterentwicklung RAG-basierter Bildungs-Chatbots am KI-Campus und benennt konkrete technische und didaktische Weiterentwicklungsbedarfe.